%%%%%%%%%%%%%%%%%%%%%%%%%%%%%%%%%%%%%%%%%%%%%%%%%%%%%%%%%%%%%%%%%%%%%%%%%%%%%%%
%%
%% Chapter 1
%%
%%%%%%%%%%%%%%%%%%%%%%%%%%%%%%%%%%%%%%%%%%%%%%%%%%%%%%%%%%%%%%%%%%%%%%%%%%%%%%%

\chapter{The Physical Problem}

The purpose of this chapter is to formulate the physical problem that
motivates the remainder of this monograph. We argue that none of the
existing approaches simultaneously provides manifest Lorentz covariance,
Hamiltonian dynamics, stochastic evolution, and a derivation of
relativistic quantum mechanics. The chapter identifies the principal
obstacles and explains why the manifestly covariant Hamiltonian dynamics
of Stueckelberg, Horwitz and Piron provides the natural deterministic
framework upon which a relativistic stochastic theory may be constructed.

%%%%%%%%%%%%%%%%%%%%%%%%%%%%%%%%%%%%%%%%%%%%%%%%%%%%%%%%%%%%%%%%%%%%%%%%%%%%%%%
\section{Chapter Overview}
%%%%%%%%%%%%%%%%%%%%%%%%%%%%%%%%%%%%%%%%%%%%%%%%%%%%%%%%%%%%%%%%%%%%%%%%%%%%%%%

The central question addressed throughout this monograph is whether
relativistic quantum mechanics can be understood as the statistical
description of an underlying stochastic dynamics formulated in a
manifestly covariant manner.

This question is considerably more subtle than its nonrelativistic
counterpart. Nelson's stochastic mechanics demonstrated that the
Schrödinger equation may be derived from a diffusion process evolving in
Newtonian time. The relativistic generalization is not obtained by a
simple replacement of Euclidean geometry by Minkowski geometry. Several
fundamental difficulties appear before one even attempts to derive a
relativistic wave equation.

The first difficulty concerns the geometry of the stochastic process.
Ordinary Brownian motion is based upon a positive-definite covariance
matrix. Minkowski space possesses an indefinite metric, so the naive
covariant generalization immediately leads to negative variances for the
spatial components of the stochastic increments.

The second difficulty concerns the evolution parameter. A stochastic
process requires an ordering parameter with respect to which random
increments are defined. Coordinate time is frame dependent, while proper
time depends upon the individual trajectory and cannot serve as a common
evolution parameter for interacting many-particle systems.

The third difficulty concerns dynamics. Conventional relativistic
mechanics is usually formulated as a constrained Hamiltonian system,
whereas stochastic dynamics is most naturally formulated as an
unconstrained evolution problem. Reconciling these two viewpoints is far
from trivial.

The manifestly covariant Hamiltonian mechanics developed by
Stueckelberg and later generalized by Horwitz and Piron resolves the
second and third difficulties by introducing an invariant evolution
parameter together with an unconstrained Hamiltonian dynamics on
space-time. The resulting formalism is structurally very close to
ordinary Hamiltonian mechanics while remaining manifestly Lorentz
covariant.

The principal thesis of this monograph is that this framework provides
the appropriate deterministic basis for a relativistic stochastic
mechanics. Rather than attempting to construct relativistic Brownian
motion directly in Minkowski space, we shall investigate stochastic
deformations of the deterministic SHP Hamiltonian flow and determine the
conditions under which the resulting statistical theory reproduces the
known equations of relativistic quantum mechanics.

This chapter has three objectives.
\begin{enumerate}
   \item
    We formulate the physical problem independently of any
    particular mathematical formalism.
   \item
    We review the principal obstacles that have prevented previous
    attempts from achieving a complete relativistic stochastic theory.
   \item
    We explain why these obstacles naturally lead to the SHP
    framework as the starting point for the developments presented in the
    subsequent chapters.
\end{enumerate}

The remainder of the monograph is organized as follows.
Chapter~2 presents a detailed review of manifestly covariant
Hamiltonian dynamics.

Chapter~3 reviews Nelson's stochastic mechanics and analyzes the
essential ingredients responsible for the emergence of the Schrödinger
equation.

Chapter~4 surveys existing relativistic stochastic processes and
identifies their strengths and limitations.

Chapter~5 reviews relativistic quantum mechanics from the viewpoint of
dynamical evolution.

Chapter~6 compares these approaches systematically and identifies the
remaining open problems.

Beginning with Chapter~7, the emphasis shifts from review to original
research. A stochastic extension of SHP dynamics is developed,
generalized to include spin, gauge interactions and many-particle
systems, and finally compared with existing formulations of relativistic
quantum theory.

%%%%%%%%%%%%%%%%%%%%%%%%%%%%%%%%%%%%%%%%%%%%%%%%%%%%%%%%%%%%%%%%%%%%%%%%%%%%%%
\section{Introduction}
%%%%%%%%%%%%%%%%%%%%%%%%%%%%%%%%%%%%%%%%%%%%%%%%%%%%%%%%%%%%%%%%%%%%%%%%%%%%%%

The relationship between classical and quantum mechanics remains one of
the central problems of theoretical physics. Although quantum mechanics
has achieved extraordinary experimental success, the origin of its
mathematical structure is still the subject of active investigation.
Among the many questions that continue to attract attention are the
physical meaning of the wave function, the origin of quantum
probabilities, the emergence of classical behavior, and the possibility
that quantum dynamics may arise from a deeper level of description.

One particularly attractive line of research is stochastic mechanics,
initiated by Fényes and developed systematically by Nelson. In this
approach the motion of an individual particle is described by a
stochastic process rather than by a deterministic trajectory. Under
appropriate assumptions, the statistical properties of this process
reproduce the Schrödinger equation. Whether or not one accepts the
interpretational aspects of stochastic mechanics, Nelson's work
demonstrated an important fact: the mathematical structure of quantum
mechanics can emerge from an underlying dynamical theory.

Despite its elegance, Nelson's construction is intrinsically
nonrelativistic. The theory is formulated in ordinary Euclidean space,
evolves with respect to Newtonian time, and relies upon concepts whose
relativistic generalization is far from straightforward. During the past
several decades numerous attempts have been made to construct a
relativistic stochastic mechanics, but no universally accepted theory
has emerged.

The difficulties are deeper than might first appear. They do not arise
because relativity and stochastic processes are individually
inconsistent. Manifestly covariant stochastic processes are known to
exist, and several mathematically rigorous constructions have been
developed. Likewise, manifestly covariant Hamiltonian dynamics has been
established through the work of Stueckelberg, Horwitz, Piron and their
collaborators. The problem is that these developments have largely
proceeded independently. A framework that simultaneously possesses
manifest Lorentz covariance, unconstrained Hamiltonian dynamics,
physically acceptable stochastic evolution, and a derivation of
relativistic quantum mechanics has not yet been achieved.

This observation suggests that the problem should be reformulated.
Rather than attempting to modify quantum mechanics directly, or to
construct relativistic Brownian motion by analogy with the
nonrelativistic theory, we seek a deterministic relativistic dynamics
that is suitable for stochastic extension. Once such a dynamics has been
identified, stochastic fluctuations may be introduced in a systematic
manner, and their physical consequences investigated.

From this viewpoint the manifestly covariant Hamiltonian dynamics of
Stueckelberg, Horwitz and Piron occupies a distinguished position.
Unlike conventional constrained formulations of relativistic mechanics,
the SHP formalism provides a genuine Hamiltonian evolution with respect
to an invariant parameter. The formal analogy with ordinary Hamiltonian
mechanics is remarkably close: canonical variables, Poisson brackets,
Hamilton equations, Liouville dynamics and Hamilton--Jacobi theory all
retain essentially their classical form. Consequently, SHP appears to be
the natural relativistic analogue of the Newtonian mechanics that serves
as the deterministic foundation of Nelson's theory.

This observation leads to the central hypothesis investigated throughout
this monograph.
The deterministic SHP dynamics plays, in relativistic physics, the same
foundational role that Newtonian mechanics plays in nonrelativistic
stochastic mechanics. A relativistic stochastic theory should therefore
be constructed not by modifying Minkowski geometry directly, but by
developing a stochastic extension of SHP Hamiltonian dynamics.

The objective of the present work is not to advocate a particular
interpretation of quantum mechanics. Instead, our aim is to investigate
whether such a stochastic extension can reproduce the known structure of
relativistic quantum mechanics while preserving the fundamental physical
principles of Lorentz covariance, Hamiltonian dynamics and gauge
invariance.

To achieve this goal we proceed in three stages.
\begin{enumerate}
   \item
    We review the existing theoretical framework, including SHP
    dynamics, Nelson's stochastic mechanics, relativistic stochastic
    processes, and relativistic quantum mechanics. Particular attention is
    paid to the assumptions underlying each approach and to the problems
    that remain unresolved.
   \item
    We perform a critical comparison of these theories in order to
    identify the precise obstacles that prevent their unification.
   \item
    Building upon this analysis, we develop a stochastic extension
    of SHP dynamics and investigate its consequences for relativistic
    quantum theory, spin, gauge interactions, and many-body systems.
\end{enumerate}
%%%%%%%%%%%%%%%%%%%%%%%%%%%%%%%%%%%%%%%%%%%%%%%%%%%%%%%%%%%%%%%%%%%%%%%%%%%%%%
\section{Requirements for a Successful Theory}
%%%%%%%%%%%%%%%%%%%%%%%%%%%%%%%%%%%%%%%%%%%%%%%%%%%%%%%%%%%%%%%%%%%%%%%%%%%%%%

Before reviewing existing approaches, it is useful to formulate the
requirements that any relativistic stochastic theory intended to
underlie quantum mechanics should satisfy. These requirements are not
postulates of a particular model; rather, they are consequences of
well-established physical principles together with the objective of
recovering relativistic quantum theory.

The purpose of this section is therefore to establish a set of criteria
against which the various existing approaches will be evaluated in the
following chapters.

%%%%%%%%%%%%%%%%%%%%%%%%%%%%%%%%%%%%%%%%%%%%%%%%%%%%%%%%%%%%%%%%%%%%%%%%%%%%%%
\subsection{Manifest Lorentz Covariance}
%%%%%%%%%%%%%%%%%%%%%%%%%%%%%%%%%%%%%%%%%%%%%%%%%%%%%%%%%%%%%%%%%%%%%%%%%%%%%%

The first and most obvious requirement is compatibility with special
relativity. The equations describing the dynamics must possess the same
form in every inertial reference frame. This requirement extends not
only to the deterministic equations of motion but also to the
statistical properties of the stochastic process.

Manifest covariance is considerably stronger than covariance of the
final observable predictions. The entire dynamical construction,
including the stochastic evolution, should be formulated in terms of
four-dimensional geometric quantities.

%%%%%%%%%%%%%%%%%%%%%%%%%%%%%%%%%%%%%%%%%%%%%%%%%%%%%%%%%%%%%%%%%%%%%%%%%%%%%%
\subsection{Hamiltonian Dynamics}
%%%%%%%%%%%%%%%%%%%%%%%%%%%%%%%%%%%%%%%%%%%%%%%%%%%%%%%%%%%%%%%%%%%%%%%%%%%%%%

The second requirement is the existence of an underlying Hamiltonian
structure.

Hamiltonian mechanics occupies a central position in modern theoretical
physics. Classical mechanics, quantum mechanics, statistical mechanics
and field theory all derive much of their mathematical structure from
Hamiltonian dynamics. It is therefore natural to require that any
stochastic extension preserve this framework as far as possible.

The deterministic limit of the theory should reproduce an unconstrained
relativistic Hamiltonian dynamics.

%%%%%%%%%%%%%%%%%%%%%%%%%%%%%%%%%%%%%%%%%%%%%%%%%%%%%%%%%%%%%%%%%%%%%%%%%%%%%%
\subsection{Invariant Evolution Parameter}
%%%%%%%%%%%%%%%%%%%%%%%%%%%%%%%%%%%%%%%%%%%%%%%%%%%%%%%%%%%%%%%%%%%%%%%%%%%%%%

A stochastic process requires an evolution parameter with respect to
which successive random increments are defined.
This parameter must satisfy three conditions.
\begin{enumerate}
\item
    It must order events uniquely along every trajectory.
\item
    Its definition must be independent of the observer.
\item
    It must remain applicable to interacting many-particle systems.
\end{enumerate}
As will be shown later in this chapter, neither coordinate time nor
proper time satisfies all three requirements simultaneously.

%%%%%%%%%%%%%%%%%%%%%%%%%%%%%%%%%%%%%%%%%%%%%%%%%%%%%%%%%%%%%%%%%%%%%%%%%%%%%%
\subsection{Positive Probability}
%%%%%%%%%%%%%%%%%%%%%%%%%%%%%%%%%%%%%%%%%%%%%%%%%%%%%%%%%%%%%%%%%%%%%%%%%%%%%%

The statistical interpretation of the theory requires the existence of a
non-negative probability density whose evolution preserves
normalization.

Although this requirement appears elementary, it is surprisingly
restrictive in relativistic theories because of the indefinite signature
of Minkowski space and the well-known difficulties associated with the
probability interpretation of relativistic wave equations.

%%%%%%%%%%%%%%%%%%%%%%%%%%%%%%%%%%%%%%%%%%%%%%%%%%%%%%%%%%%%%%%%%%%%%%%%%%%%%%
\subsection{Correct Nonrelativistic Limit}
%%%%%%%%%%%%%%%%%%%%%%%%%%%%%%%%%%%%%%%%%%%%%%%%%%%%%%%%%%%%%%%%%%%%%%%%%%%%%%

Any acceptable relativistic theory must reduce continuously to the
established nonrelativistic theory in the limit
\[v/c \rightarrow 0.\]
In particular, the stochastic dynamics should reproduce Nelson's
construction whenever relativistic effects become negligible.

This requirement provides one of the most important consistency checks
for any proposed relativistic extension.

%%%%%%%%%%%%%%%%%%%%%%%%%%%%%%%%%%%%%%%%%%%%%%%%%%%%%%%%%%%%%%%%%%%%%%%%%%%%%%
\subsection{Gauge Invariance}
%%%%%%%%%%%%%%%%%%%%%%%%%%%%%%%%%%%%%%%%%%%%%%%%%%%%%%%%%%%%%%%%%%%%%%%%%%%%%%

Interactions with external electromagnetic fields should be introduced
through a manifestly gauge-covariant prescription. The stochastic
dynamics must preserve gauge invariance in exactly the same sense as the
underlying deterministic Hamiltonian theory.

%%%%%%%%%%%%%%%%%%%%%%%%%%%%%%%%%%%%%%%%%%%%%%%%%%%%%%%%%%%%%%%%%%%%%%%%%%%%%%
\subsection{Spin}
%%%%%%%%%%%%%%%%%%%%%%%%%%%%%%%%%%%%%%%%%%%%%%%%%%%%%%%%%%%%%%%%%%%%%%%%%%%%%%

A complete relativistic theory must ultimately describe particles with
spin.

Whether spin should emerge from the stochastic dynamics or be introduced
through additional dynamical variables remains an open question.
Regardless of the answer, the formalism should admit a natural extension
to particles of arbitrary spin.

%%%%%%%%%%%%%%%%%%%%%%%%%%%%%%%%%%%%%%%%%%%%%%%%%%%%%%%%%%%%%%%%%%%%%%%%%%%%%%
\subsection{Many-Particle Systems}
%%%%%%%%%%%%%%%%%%%%%%%%%%%%%%%%%%%%%%%%%%%%%%%%%%%%%%%%%%%%%%%%%%%%%%%%%%%%%%

Finally, the theory should extend naturally to systems containing many
interacting particles.

The existence of a common evolution parameter is particularly important
in this context. Without such a parameter the formulation of interacting
relativistic stochastic dynamics becomes considerably more complicated.

%%%%%%%%%%%%%%%%%%%%%%%%%%%%%%%%%%%%%%%%%%%%%%%%%%%%%%%%%%%%%%%%%%%%%%%%%%%%%%
\subsection{Summary of the Requirements}
%%%%%%%%%%%%%%%%%%%%%%%%%%%%%%%%%%%%%%%%%%%%%%%%%%%%%%%%%%%%%%%%%%%%%%%%%%%%%%

The preceding discussion identifies the principal criteria that will be
used throughout this monograph.
A successful relativistic stochastic theory should simultaneously
possess
\begin{enumerate}
\item manifest Lorentz covariance,
\item unconstrained Hamiltonian dynamics,
\item an invariant evolution parameter,
\item a positive probabilistic interpretation,
\item gauge invariance,
\item a correct nonrelativistic limit,
\item a natural treatment of spin,
\item a consistent many-particle formulation.
\end{enumerate}
No presently available theory satisfies all of these requirements
simultaneously. The remainder of Part~I is devoted to examining how
close the existing approaches come to this objective and to identifying
the obstacles that remain.
%%%%%%%%%%%%%%%%%%%%%%%%%%%%%%%%%%%%%%%%%%%%%%%%%%%%%%%%%%%%%%%%%%%%%%%%%%%%%%
\section{Failure of Na\"ive Relativistic Brownian Motion}
%%%%%%%%%%%%%%%%%%%%%%%%%%%%%%%%%%%%%%%%%%%%%%%%%%%%%%%%%%%%%%%%%%%%%%%%%%%%%%

The construction of a relativistic stochastic mechanics might appear, at
first sight, to be straightforward. Since Brownian motion is well
understood in Euclidean space, one could simply attempt to replace the
Euclidean metric by the Minkowski metric while preserving the remaining
structure of the theory.

This seemingly natural procedure fails immediately.

The failure is not a technical complication but a fundamental
consequence of the indefinite geometry of Minkowski space. The purpose of
this section is to demonstrate this fact explicitly.

%%%%%%%%%%%%%%%%%%%%%%%%%%%%%%%%%%%%%%%%%%%%%%%%%%%%%%%%%%%%%%%%%%%%%%%%%%%%%%
\subsection{Brownian Motion in Euclidean Space}
%%%%%%%%%%%%%%%%%%%%%%%%%%%%%%%%%%%%%%%%%%%%%%%%%%%%%%%%%%%%%%%%%%%%%%%%%%%%%%

Consider first the standard Wiener process in three-dimensional
Euclidean space.
The infinitesimal displacement is written as
\[dX^i=b^i\,dt+dW^i ,\]
where \(b^i\) denotes the drift velocity and \(dW^i\) is the random
increment.

The defining statistical properties of the Wiener process are
\[\langle dW^i\rangle =0,\]
and
\[\langle dW^i dW^j\rangle=2D\,\delta^{ij}\,dt ,\]
where \(D>0\) is the diffusion coefficient.

The covariance matrix
\[C^{ij}=2D\,\delta^{ij}\,dt\]
is positive definite.
Consequently,
\[\langle (dW^i)^2\rangle=2D\,dt>0,\]
for every spatial component.
This positivity is essential. It guarantees that the stochastic
increments possess physically meaningful variances and allows the
construction of the associated Fokker--Planck equation.

%%%%%%%%%%%%%%%%%%%%%%%%%%%%%%%%%%%%%%%%%%%%%%%%%%%%%%%%%%%%%%%%%%%%%%%%%%%%%%
\subsection{A Na\"ive Covariant Generalization}
%%%%%%%%%%%%%%%%%%%%%%%%%%%%%%%%%%%%%%%%%%%%%%%%%%%%%%%%%%%%%%%%%%%%%%%%%%%%%%

Suppose one attempts to formulate an analogous stochastic process in
Minkowski space.
The natural covariant generalization appears to be
\[dX^\mu=b^\mu d\tau+dW^\mu ,\]
where \(\tau\) is an invariant evolution parameter.
Assuming Lorentz covariance, one would naturally postulate
\[\langle dW^\mu\rangle =0,\]
together with
\[\boxed{\langle dW^\mu dW^\nu\rangle=2D\,\eta^{\mu\nu}\,d\tau .}\]
At first glance this expression possesses every desired symmetry.
It is manifestly Lorentz covariant.
It reduces to the ordinary Wiener process if the metric is replaced by
its Euclidean counterpart.
Unfortunately, it cannot describe a physical stochastic process.

%%%%%%%%%%%%%%%%%%%%%%%%%%%%%%%%%%%%%%%%%%%%%%%%%%%%%%%%%%%%%%%%%%%%%%%%%%%%%%
\subsection{The Problem of Negative Variances}
%%%%%%%%%%%%%%%%%%%%%%%%%%%%%%%%%%%%%%%%%%%%%%%%%%%%%%%%%%%%%%%%%%%%%%%%%%%%%%

Throughout this monograph we adopt the metric convention
\[\eta_{\mu\nu}=\mathrm{diag}(+,-,-,-).\]
Equation (1.25) therefore implies
\[\langle(dW^0)^2\rangle=2D\,d\tau ,\]
whereas
\[\langle(dW^1)^2\rangle=-2D\,d\tau ,\]
and similarly
\[\langle(dW^2)^2\rangle=-2D\,d\tau ,\]
\[\langle(dW^3)^2\rangle=-2D\,d\tau .\]
The spatial variances are therefore negative.

Since the variance of a real random variable satisfies
\[\mathrm{Var}(X)=\langle X^2\rangle-\langle X\rangle^2\ge 0 ,\]
the covariance matrix above cannot correspond to any ordinary stochastic
process.
The difficulty arises solely from the Lorentzian signature of space-time
and is therefore independent of the particular physical model.

%%%%%%%%%%%%%%%%%%%%%%%%%%%%%%%%%%%%%%%%%%%%%%%%%%%%%%%%%%%%%%%%%%%%%%%%%%%%%%
\subsection{Why the Problem is Fundamental}
%%%%%%%%%%%%%%%%%%%%%%%%%%%%%%%%%%%%%%%%%%%%%%%%%%%%%%%%%%%%%%%%%%%%%%%%%%%%%%

It is tempting to regard the appearance of negative variances as merely
a technical inconvenience that could be eliminated by modifying the
diffusion coefficient or by introducing an anisotropic covariance
matrix.
Such modifications do not resolve the underlying difficulty.

The covariance tensor of a stochastic process must define a positive
semi-definite quadratic form,
\[A_\mu\langle dW^\mu dW^\nu\rangle A_\nu\ge 0\]
for every real four-vector \(A^\mu\).
By contrast, the Minkowski metric satisfies
\[A_\mu\eta^{\mu\nu}A_\nu=A_0^2-\mathbf{A}^2 ,\]
which assumes both positive and negative values depending upon the
direction of \(A^\mu\).
Therefore the Minkowski metric itself cannot serve as the covariance
tensor of an ordinary diffusion process.
This conclusion is independent of the value of the diffusion
coefficient, the choice of drift velocity, or the particular form of the
Hamiltonian.

%%%%%%%%%%%%%%%%%%%%%%%%%%%%%%%%%%%%%%%%%%%%%%%%%%%%%%%%%%%%%%%%%%%%%%%%%%%%%%
\subsection{Possible Ways Around the Difficulty}
%%%%%%%%%%%%%%%%%%%%%%%%%%%%%%%%%%%%%%%%%%%%%%%%%%%%%%%%%%%%%%%%%%%%%%%%%%%%%%

The failure of the naive construction does not imply that relativistic
stochastic mechanics is impossible.
Rather, it demonstrates that at least one of the assumptions underlying
the nonrelativistic Wiener process must be abandoned or generalized.
Several possibilities have been explored in the literature.

One may restrict the stochastic process to timelike trajectories, define
the diffusion on the mass shell or on the unit hyperboloid of
four-velocities, introduce stochastic evolution in proper time or in an
independent invariant parameter, or replace the ordinary Wiener process
by a more general Markov process on a Lorentzian manifold.

These approaches will be examined in detail in Chapter~4.
At present we merely emphasize that the problem appears before any
specific stochastic model has been chosen. It is therefore a geometric
obstacle rather than a dynamical one.

%%%%%%%%%%%%%%%%%%%%%%%%%%%%%%%%%%%%%%%%%%%%%%%%%%%%%%%%%%%%%%%%%%%%%%%%%%%%%%
\subsection{Conclusions}
%%%%%%%%%%%%%%%%%%%%%%%%%%%%%%%%%%%%%%%%%%%%%%%%%%%%%%%%%%%%%%%%%%%%%%%%%%%%%%

The preceding analysis establishes the first fundamental obstacle to a
relativistic stochastic mechanics.
The naive replacement
\[\delta^{ij}\longrightarrow\eta^{\mu\nu}\]
does not produce a physically meaningful relativistic Wiener process.
Consequently, relativistic stochastic mechanics requires substantially
more than a covariant rewriting of the nonrelativistic theory.

The next question is therefore not how to define the covariance matrix,
but rather how to formulate the evolution itself.
As we shall see in the following section, this immediately leads to the
problem of identifying an appropriate invariant evolution parameter.
%%%%%%%%%%%%%%%%%%%%%%%%%%%%%%%%%%%%%%%%%%%%%%%%%%%%%%%%%%%%%%%%%%%%%%%%%%%%%%
\section{The Evolution Parameter Problem}
%%%%%%%%%%%%%%%%%%%%%%%%%%%%%%%%%%%%%%%%%%%%%%%%%%%%%%%%%%%%%%%%%%%%%%%%%%%%%%

The failure of the naive relativistic generalization of Brownian motion
raises a more fundamental question. A stochastic process is defined by a
sequence of infinitesimal random increments,
\[X \longrightarrow X+dX,\]
ordered by an evolution parameter. Before constructing a relativistic
stochastic theory one must therefore answer a simpler question:
\begin{center}
\emph{What is the appropriate evolution parameter for relativistic
dynamics?}
\end{center}

In nonrelativistic mechanics the answer is immediate. Every trajectory
is parameterized by the universal Newtonian time,
\[\mathbf{x}=\mathbf{x}(t),\]
and the stochastic process
\[\mathbf{X}(t)\]
is defined with respect to this common temporal parameter.

Relativity changes this situation profoundly. The absence of universal
simultaneity implies that the notion of time itself depends upon the
observer. Consequently, the construction of a relativistic stochastic
process requires a careful examination of possible evolution
parameters.

%%%%%%%%%%%%%%%%%%%%%%%%%%%%%%%%%%%%%%%%%%%%%%%%%%%%%%%%%%%%%%%%%%%%%%%%%%%%%%
\subsection{Coordinate Time}
%%%%%%%%%%%%%%%%%%%%%%%%%%%%%%%%%%%%%%%%%%%%%%%%%%%%%%%%%%%%%%%%%%%%%%%%%%%%%%
The most obvious choice is the coordinate time,
\[x^0=ct.\]
This parameter possesses a clear operational interpretation within a
given inertial reference frame.
However, coordinate time is not Lorentz invariant.
Under a Lorentz transformation,
\[x'^0=\gamma\left(x^0-\beta x^1\right),\]
where
\[\gamma=\frac{1}{\sqrt{1-\beta^2}}.\]

Different inertial observers therefore assign different temporal
coordinates to the same event.
A stochastic process defined by
\[X^\mu(t)\]
is consequently tied to a particular Lorentz frame. Although the final
observable predictions may still be Lorentz invariant, manifest
covariance of the dynamical construction is lost.

Since one of the principal objectives of the present work is to preserve
manifest covariance throughout, coordinate time is not an appropriate
choice.

%%%%%%%%%%%%%%%%%%%%%%%%%%%%%%%%%%%%%%%%%%%%%%%%%%%%%%%%%%%%%%%%%%%%%%%%%%%%%%
\subsection{Proper Time}
%%%%%%%%%%%%%%%%%%%%%%%%%%%%%%%%%%%%%%%%%%%%%%%%%%%%%%%%%%%%%%%%%%%%%%%%%%%%%%
The next possibility is the invariant proper time,
\[ds^2=dx^\mu dx_\mu.\]
Proper time possesses several attractive properties.
It is Lorentz invariant.
It has direct physical meaning.
For free particles it parameterizes the worldline uniquely.
Nevertheless, proper time is not suitable as a universal evolution
parameter.
\begin{enumerate}
   \item
    Proper time exists only along timelike trajectories. It cannot
    parameterize null worldlines and becomes problematic in processes
    involving virtual or off-shell particles.
   \item
    Proper time depends upon the trajectory itself. Two interacting
    particles generally possess different proper times,
    \[s_1 \neq s_2,\]
    making it impossible to formulate a common dynamical evolution for an
    interacting many-body system.
   \item
    Proper time is not an independent dynamical variable. It is
    determined by the motion rather than specifying it.
\end{enumerate}
These difficulties severely limit its usefulness as the evolution
parameter of a stochastic dynamics.
%%%%%%%%%%%%%%%%%%%%%%%%%%%%%%%%%%%%%%%%%%%%%%%%%%%%%%%%%%%%%%%%%%%%%%%%%%%%%%
\subsection{Affine Parameters}
%%%%%%%%%%%%%%%%%%%%%%%%%%%%%%%%%%%%%%%%%%%%%%%%%%%%%%%%%%%%%%%%%%%%%%%%%%%%%%
One may introduce an arbitrary affine parameter,
\[\lambda,\]
and write
\[x^\mu=x^\mu(\lambda).\]
Such a parameter avoids the difficulties associated with coordinate time
and proper time.
However, its physical interpretation remains ambiguous.

Unless additional dynamical structure is introduced, different choices
of affine parameter merely correspond to different parametrizations of
the same trajectory.
A stochastic process constructed with respect to an arbitrary parameter
therefore lacks an intrinsic physical meaning.
%%%%%%%%%%%%%%%%%%%%%%%%%%%%%%%%%%%%%%%%%%%%%%%%%%%%%%%%%%%%%%%%%%%%%%%%%%%%%%
\subsection{The Requirements}
%%%%%%%%%%%%%%%%%%%%%%%%%%%%%%%%%%%%%%%%%%%%%%%%%%%%%%%%%%%%%%%%%%%%%%%%%%%%%%
The preceding discussion suggests that the evolution parameter required
for relativistic stochastic dynamics must satisfy several conditions.
It should
\begin{enumerate}
\item be Lorentz invariant,
\item exist independently of the particular trajectory,
\item provide a unique ordering of events,
\item remain applicable to interacting many-particle systems,
\item generate an unconstrained Hamiltonian evolution.
\end{enumerate}
Remarkably, no parameter occurring naturally in conventional
relativistic mechanics satisfies all of these requirements
simultaneously.
%%%%%%%%%%%%%%%%%%%%%%%%%%%%%%%%%%%%%%%%%%%%%%%%%%%%%%%%%%%%%%%%%%%%%%%%%%%%%%
\subsection{An Invariant Dynamical Parameter}
%%%%%%%%%%%%%%%%%%%%%%%%%%%%%%%%%%%%%%%%%%%%%%%%%%%%%%%%%%%%%%%%%%%%%%%%%%%%%%
The difficulty described above motivated the introduction of an
independent invariant evolution parameter by Stueckelberg and,
subsequently, its systematic development by Horwitz and Piron.

Instead of regarding the temporal coordinate as the evolution
parameter, one introduces an independent scalar quantity,
\[\tau,\]
and regards all four space-time coordinates as dynamical variables,
\[x^\mu=x^\mu(\tau).\]
Unlike coordinate time, the parameter \(\tau\) is Lorentz invariant.
Unlike proper time, it is independent of the particular trajectory.
Unlike an arbitrary affine parameter, it acquires a precise dynamical
meaning through Hamilton's equations of motion.
Consequently, \(\tau\) satisfies exactly the requirements identified
above.
%%%%%%%%%%%%%%%%%%%%%%%%%%%%%%%%%%%%%%%%%%%%%%%%%%%%%%%%%%%%%%%%%%%%%%%%%%%%%%
\subsection{Physical Interpretation}
%%%%%%%%%%%%%%%%%%%%%%%%%%%%%%%%%%%%%%%%%%%%%%%%%%%%%%%%%%%%%%%%%%%%%%%%%%%%%%
The introduction of the invariant parameter \(\tau\) represents a
conceptual change in the formulation of relativistic dynamics.
Coordinate time,
\[x^0,\]
is promoted to the status of an ordinary dynamical variable on an equal
footing with the spatial coordinates,
\[x^1,\;x^2,\;x^3.\]
The evolution of the entire four-dimensional trajectory is generated
with respect to the invariant parameter \(\tau\).

From the viewpoint of deterministic mechanics this modification appears
minor.
From the viewpoint of stochastic dynamics it is profound.
A stochastic process requires precisely such an independent ordering
parameter. The introduction of \(\tau\) therefore removes one of the
principal conceptual obstacles that prevented a manifestly covariant
generalization of Nelson's construction.
%%%%%%%%%%%%%%%%%%%%%%%%%%%%%%%%%%%%%%%%%%%%%%%%%%%%%%%%%%%%%%%%%%%%%%%%%%%%%%
\subsection{Conclusions}
%%%%%%%%%%%%%%%%%%%%%%%%%%%%%%%%%%%%%%%%%%%%%%%%%%%%%%%%%%%%%%%%%%%%%%%%%%%%%%
The principal difficulty of relativistic stochastic mechanics is not the
construction of a Lorentz-covariant diffusion process itself.
Rather, it is the absence, within conventional relativistic mechanics,
of a physically meaningful evolution parameter possessing the properties
required by stochastic dynamics.
The introduction of the invariant parameter \(\tau\) resolves this
difficulty and naturally leads to the manifestly covariant Hamiltonian
formalism of Stueckelberg, Horwitz and Piron.
The structure of this deterministic theory forms the subject of the next
section.
%%%%%%%%%%%%%%%%%%%%%%%%%%%%%%%%%%%%%%%%%%%%%%%%%%%%%%%%%%%%%%%%%%%%%%%%%%%%%%
\section{Manifestly Covariant Hamiltonian Dynamics}
%%%%%%%%%%%%%%%%%%%%%%%%%%%%%%%%%%%%%%%%%%%%%%%%%%%%%%%%%%%%%%%%%%%%%%%%%%%%%%
The preceding discussion identified the principal obstacle to the
construction of a relativistic stochastic mechanics. A stochastic process
requires an invariant evolution parameter together with a deterministic
dynamical law governing the evolution of the system. Conventional
relativistic mechanics provides neither in a form suitable for an
unconstrained stochastic theory.

The purpose of this section is to introduce a deterministic relativistic
dynamics satisfying the requirements established above. A detailed
development will be presented in Chapter~2. Here we discuss only the
essential physical ideas.
%%%%%%%%%%%%%%%%%%%%%%%%%%%%%%%%%%%%%%%%%%%%%%%%%%%%%%%%%%%%%%%%%%%%%%%%%%%%%%
\subsection{From Constrained to Unconstrained Dynamics}
%%%%%%%%%%%%%%%%%%%%%%%%%%%%%%%%%%%%%%%%%%%%%%%%%%%%%%%%%%%%%%%%%%%%%%%%%%%%%%
The standard formulation of relativistic particle mechanics is based upon
the mass-shell constraint
\[p^\mu p_\mu = m^2c^2.\]
The particle trajectory is usually parametrized by the proper time or by
an arbitrary affine parameter, while the dynamics is generated through
the constraint rather than by an ordinary Hamiltonian.
This approach is entirely adequate for many classical problems.
Nevertheless, it differs fundamentally from the Hamiltonian mechanics
used in nonrelativistic physics.

In particular,
\begin{itemize}
\item the evolution parameter is not independent,
\item the Hamiltonian is replaced by a constraint,
\item the temporal coordinate is not treated as an ordinary canonical
variable.
\end{itemize}
These features complicate the construction of a stochastic extension.
%%%%%%%%%%%%%%%%%%%%%%%%%%%%%%%%%%%%%%%%%%%%%%%%%%%%%%%%%%%%%%%%%%%%%%%%%%%%%%
\subsection{The SHP Hamiltonian}
%%%%%%%%%%%%%%%%%%%%%%%%%%%%%%%%%%%%%%%%%%%%%%%%%%%%%%%%%%%%%%%%%%%%%%%%%%%%%%
Stueckelberg proposed an alternative formulation in which the evolution
is generated with respect to an invariant parameter,
\[\tau,\]
independent of the space-time coordinates.
The canonical variables are
\[x^\mu(\tau),\qquad p_\mu(\tau),\]
and the dynamics is generated by a Hamiltonian function
\[K(x,p).\]
For a free particle,
\[\boxed{K_0=\frac{p^\mu p_\mu}{2M},}\]
where \(M\) is a constant parameter with the dimensions of mass.

Interactions are introduced through an invariant scalar potential,
\[K=\frac{p^\mu p_\mu}{2M}+V(x).\]
The significance of this expression lies not in its particular form but
in the fact that it generates an unconstrained Hamiltonian evolution.
%%%%%%%%%%%%%%%%%%%%%%%%%%%%%%%%%%%%%%%%%%%%%%%%%%%%%%%%%%%%%%%%%%%%%%%%%%%%%%
\subsection{Hamilton's Equations}
%%%%%%%%%%%%%%%%%%%%%%%%%%%%%%%%%%%%%%%%%%%%%%%%%%%%%%%%%%%%%%%%%%%%%%%%%%%%%%
The equations of motion retain exactly the canonical form familiar from
ordinary Hamiltonian mechanics,
\[\boxed{\frac{dx^\mu}{d\tau}=\frac{\partial K}{\partial p_\mu},}\]
\[\boxed{\frac{dp_\mu}{d\tau}=-\frac{\partial K}{\partial x^\mu}.}\]

For the free Hamiltonian these equations become
\[\frac{dx^\mu}{d\tau}=\frac{p^\mu}{M},\]
\[\frac{dp_\mu}{d\tau}=0.\]
Thus the free particle follows a straight worldline in space-time, while
all four coordinates evolve dynamically with respect to the invariant
parameter.
%%%%%%%%%%%%%%%%%%%%%%%%%%%%%%%%%%%%%%%%%%%%%%%%%%%%%%%%%%%%%%%%%%%%%%%%%%%%%%
\subsection{Off-Shell Dynamics}
%%%%%%%%%%%%%%%%%%%%%%%%%%%%%%%%%%%%%%%%%%%%%%%%%%%%%%%%%%%%%%%%%%%%%%%%%%%%%%
One of the most distinctive features of the SHP formalism is that the
quantity
\[p^\mu p_\mu\]
is not imposed as a fixed constraint.
Instead, it becomes a dynamical quantity whose value may evolve under the
influence of interactions.
The invariant mass,
\[m^2=p^\mu p_\mu ,\]
is therefore not necessarily constant during the evolution.
Only in the free-particle limit does one recover
\[m=M,\]
so that the conventional mass-shell condition is restored.

This off-shell character plays an essential role in many applications of
the SHP theory and may prove equally important in the construction of a
stochastic extension.
%%%%%%%%%%%%%%%%%%%%%%%%%%%%%%%%%%%%%%%%%%%%%%%%%%%%%%%%%%%%%%%%%%%%%%%%%%%%%%
\subsection{Why SHP is Different}
%%%%%%%%%%%%%%%%%%%%%%%%%%%%%%%%%%%%%%%%%%%%%%%%%%%%%%%%%%%%%%%%%%%%%%%%%%%%%%
At first sight the introduction of an invariant parameter might appear to
be merely a reparametrization of conventional relativistic mechanics.
This impression is incorrect.
The SHP formalism is characterized by several features absent from the
standard constrained theory.
\begin{enumerate}
   \item
    The evolution parameter is independent of the trajectory.
   \item
    Hamiltonian evolution is generated by an ordinary scalar
    Hamiltonian rather than by a constraint.
   \item
    The temporal coordinate is treated as a canonical dynamical
    variable on the same footing as the spatial coordinates.
   \item
    The formalism naturally accommodates interacting many-particle
    systems through a common invariant evolution parameter.
\end{enumerate}
These properties make SHP structurally much closer to Newtonian
Hamiltonian mechanics than to conventional constrained relativistic
dynamics.
%%%%%%%%%%%%%%%%%%%%%%%%%%%%%%%%%%%%%%%%%%%%%%%%%%%%%%%%%%%%%%%%%%%%%%%%%%%%%%
\subsection{Connection with Stochastic Mechanics}
%%%%%%%%%%%%%%%%%%%%%%%%%%%%%%%%%%%%%%%%%%%%%%%%%%%%%%%%%%%%%%%%%%%%%%%%%%%%%%
The importance of the SHP formalism for the present work does not lie
solely in its manifest covariance.
Equally important is its Hamiltonian structure.
Nelson's stochastic mechanics begins with deterministic Newtonian
dynamics and replaces deterministic trajectories by stochastic ones.
The same philosophy can now be adopted in the relativistic domain.

Instead of attempting to construct stochastic motion directly in
Minkowski space, we begin with the deterministic Hamiltonian flow
generated by the SHP Hamiltonian and investigate stochastic
deformations of this flow.
From this viewpoint the SHP formalism plays the same conceptual role in
relativistic stochastic mechanics that Newtonian mechanics plays in
Nelson's original construction.
%%%%%%%%%%%%%%%%%%%%%%%%%%%%%%%%%%%%%%%%%%%%%%%%%%%%%%%%%%%%%%%%%%%%%%%%%%%%%%
\subsection{Conclusions}
%%%%%%%%%%%%%%%%%%%%%%%%%%%%%%%%%%%%%%%%%%%%%%%%%%%%%%%%%%%%%%%%%%%%%%%%%%%%%%
The SHP formalism provides a manifestly covariant Hamiltonian dynamics
possessing precisely those structural features required for a
relativistic stochastic theory.
It does not by itself introduce stochasticity.
Rather, it supplies the deterministic framework upon which such a theory
may be constructed.
The next section discusses the general principles by which deterministic
Hamiltonian dynamics may be extended to include stochastic evolution.
%%%%%%%%%%%%%%%%%%%%%%%%%%%%%%%%%%%%%%%%%%%%%%%%%%%%%%%%%%%%%%%%%%%%%%%%%%%%%%
\section{From Deterministic to Stochastic Dynamics}
%%%%%%%%%%%%%%%%%%%%%%%%%%%%%%%%%%%%%%%%%%%%%%%%%%%%%%%%%%%%%%%%%%%%%%%%%%%%%%
The preceding sections have identified the deterministic framework that
will serve as the starting point for the remainder of this monograph.
The next question is how stochastic effects should be incorporated into
this dynamics.
It is important to emphasize that the problem is not the construction of
an arbitrary relativistic stochastic process. The objective is much more
specific. We seek a stochastic generalization of an already existing
Hamiltonian dynamics.
This distinction is fundamental.

In classical statistical mechanics the equations of Newton are not
replaced by stochastic equations. Rather, stochastic forces are
introduced as corrections to an underlying deterministic dynamics.
Similarly, the Langevin equation supplements Newton's second law by
adding fluctuating and dissipative forces while preserving the
deterministic motion as the mean evolution.
The same philosophy will be adopted here.
The deterministic SHP equations are regarded as the zeroth-order
approximation describing the average motion of the particle. Random
fluctuations are then introduced as additional dynamical contributions.
%%%%%%%%%%%%%%%%%%%%%%%%%%%%%%%%%%%%%%%%%%%%%%%%%%%%%%%%%%%%%%%%%%%%%%%%%%%%%%
\subsection{A General Stochastic Extension}
%%%%%%%%%%%%%%%%%%%%%%%%%%%%%%%%%%%%%%%%%%%%%%%%%%%%%%%%%%%%%%%%%%%%%%%%%%%%%%
Let
\[K(x,p)\]
be the Hamiltonian generating the deterministic dynamics.
The canonical equations are
\[\frac{dx^\mu}{d\tau}=\frac{\partial K}{\partial p_\mu},\]
\[\frac{dp_\mu}{d\tau}=-\frac{\partial K}{\partial x^\mu}.\]
A stochastic extension should reduce to these equations whenever the
random fluctuations vanish.

The most general infinitesimal evolution may therefore be written in the
schematic form
\[dx^\mu=\frac{\partial K}{\partial p_\mu}\,d\tau+d\xi^\mu,\]
\[dp_\mu=-\frac{\partial K}{\partial x^\mu}\,d\tau+d\eta_\mu,\]
where
\[d\xi^\mu,\qquad d\eta_\mu,\]
denote stochastic increments.
At this stage no assumptions are made concerning their statistical
properties.
%%%%%%%%%%%%%%%%%%%%%%%%%%%%%%%%%%%%%%%%%%%%%%%%%%%%%%%%%%%%%%%%%%%%%%%%%%%%%%
\subsection{Deterministic Limit}
%%%%%%%%%%%%%%%%%%%%%%%%%%%%%%%%%%%%%%%%%%%%%%%%%%%%%%%%%%%%%%%%%%%%%%%%%%%%%%
Consistency requires that the deterministic equations be recovered when
the stochastic fluctuations disappear.
If
\[d\xi^\mu=0,\qquad d\eta_\mu=0,\]
the stochastic equations reduce immediately to the canonical SHP
equations.

More generally, if the fluctuations possess vanishing expectation
values,
\[\langle d\xi^\mu\rangle =0,\]
\[\langle d\eta_\mu\rangle =0,\]
then the ensemble averages satisfy
\[\left\langle dx^\mu\right\rangle=\frac{\partial K}{\partial p_\mu}\,d\tau,\]
\[\left\langle dp_\mu\right\rangle=-\frac{\partial K}{\partial x^\mu}\,d\tau.\]
Thus the deterministic SHP motion represents the mean evolution of the
stochastic system.
%%%%%%%%%%%%%%%%%%%%%%%%%%%%%%%%%%%%%%%%%%%%%%%%%%%%%%%%%%%%%%%%%%%%%%%%%%%%%%
\subsection{What Remains Unknown}
%%%%%%%%%%%%%%%%%%%%%%%%%%%%%%%%%%%%%%%%%%%%%%%%%%%%%%%%%%%%%%%%%%%%%%%%%%%%%%
Equations (1.xx) and (1.xx) are intentionally incomplete.
The stochastic increments have not yet been specified.
Several fundamental questions remain unanswered.
Are the increments Markovian?
Do they possess independent increments?
Is the diffusion Gaussian?
Is the covariance tensor local?
Can the fluctuations be represented by a Wiener process, or is a more
general stochastic process required?
These questions cannot be answered on purely mathematical grounds.
Their resolution must be guided by physical principles together with the
requirement that the resulting statistical theory reproduce relativistic
quantum mechanics.
%%%%%%%%%%%%%%%%%%%%%%%%%%%%%%%%%%%%%%%%%%%%%%%%%%%%%%%%%%%%%%%%%%%%%%%%%%%%%%
\subsection{Comparison with Nelson's Construction}
%%%%%%%%%%%%%%%%%%%%%%%%%%%%%%%%%%%%%%%%%%%%%%%%%%%%%%%%%%%%%%%%%%%%%%%%%%%%%%
The strategy adopted here differs subtly from the historical
development of stochastic mechanics.
Nelson began with a diffusion process in Euclidean configuration space
and showed that, under suitable assumptions, the Schrödinger equation
follows.

The present work reverses the logical order.
We first identify the deterministic relativistic Hamiltonian dynamics.
Only afterwards do we investigate admissible stochastic modifications of
that dynamics.
This distinction reflects the greater complexity of relativistic
kinematics and dynamics.
In the nonrelativistic theory the deterministic background is uniquely
provided by Newtonian mechanics.
In relativistic physics the choice of deterministic framework is itself
a nontrivial question.
%%%%%%%%%%%%%%%%%%%%%%%%%%%%%%%%%%%%%%%%%%%%%%%%%%%%%%%%%%%%%%%%%%%%%%%%%%%%%%
\subsection{Criteria for the Stochastic Process}
%%%%%%%%%%%%%%%%%%%%%%%%%%%%%%%%%%%%%%%%%%%%%%%%%%%%%%%%%%%%%%%%%%%%%%%%%%%%%%
Although the detailed mathematical structure of the stochastic process
will be developed only in later chapters, the preceding discussion
already imposes several important requirements.

The stochastic evolution should
\begin{enumerate}
\item preserve manifest Lorentz covariance,
\item reduce continuously to deterministic SHP dynamics,
\item possess a well-defined probabilistic interpretation,
\item admit a consistent transport equation for probability densities,
\item reproduce Nelson's stochastic mechanics in the
nonrelativistic limit,
\item lead ultimately to the experimentally verified equations of
relativistic quantum mechanics.
\end{enumerate}
These requirements provide the guiding principles for the construction
developed in Part~III of this monograph.
%%%%%%%%%%%%%%%%%%%%%%%%%%%%%%%%%%%%%%%%%%%%%%%%%%%%%%%%%%%%%%%%%%%%%%%%%%%%%%
\subsection{Conclusions}
%%%%%%%%%%%%%%%%%%%%%%%%%%%%%%%%%%%%%%%%%%%%%%%%%%%%%%%%%%%%%%%%%%%%%%%%%%%%%%
The introduction of stochasticity should be regarded as a deformation of
deterministic Hamiltonian evolution rather than as an independent
kinematical postulate.
This viewpoint preserves the geometric and dynamical structure of the
underlying SHP theory while allowing additional statistical effects to
be incorporated in a systematic manner.
The next two chapters review the principal ingredients required for such
a construction: Nelson's stochastic mechanics and the existing theories
of relativistic stochastic processes.
%%%%%%%%%%%%%%%%%%%%%%%%%%%%%%%%%%%%%%%%%%%%%%%%%%%%%%%%%%%%%%%%%%%%%%%%%%%%%%
\section{Existing Research Programmes}
%%%%%%%%%%%%%%%%%%%%%%%%%%%%%%%%%%%%%%%%%%%%%%%%%%%%%%%%%%%%%%%%%%%%%%%%%%%%%%
The problem of constructing a relativistic stochastic mechanics has
attracted sustained attention for more than half a century. Numerous
approaches have been proposed, each motivated by different physical and
mathematical considerations. Despite significant progress, no generally
accepted theory has yet emerged.

Rather than reviewing the historical development chronologically, it is
more instructive to classify the existing approaches according to the
physical problems they attempt to solve. This viewpoint will also
clarify the role of the subsequent chapters of the present monograph.
%%%%%%%%%%%%%%%%%%%%%%%%%%%%%%%%%%%%%%%%%%%%%%%%%%%%%%%%%%%%%%%%%%%%%%%%%%%%%%
\subsection{Stochastic Mechanics}
%%%%%%%%%%%%%%%%%%%%%%%%%%%%%%%%%%%%%%%%%%%%%%%%%%%%%%%%%%%%%%%%%%%%%%%%%%%%%%
The first research programme originates in the work of Fényes and
Nelson. Its objective is to derive the Schrödinger equation from an
underlying stochastic dynamics without assuming the existence of a wave
function a priori.

The principal achievement of this programme is the demonstration that
quantum dynamics may emerge from a diffusion process formulated in
configuration space.
Its principal limitation is that the construction relies upon absolute
Newtonian time and Euclidean geometry. A straightforward relativistic
generalization is therefore impossible.
A detailed review of this programme will be presented in Chapter~3.
%%%%%%%%%%%%%%%%%%%%%%%%%%%%%%%%%%%%%%%%%%%%%%%%%%%%%%%%%%%%%%%%%%%%%%%%%%%%%%
\subsection{Manifestly Covariant Hamiltonian Dynamics}
%%%%%%%%%%%%%%%%%%%%%%%%%%%%%%%%%%%%%%%%%%%%%%%%%%%%%%%%%%%%%%%%%%%%%%%%%%%%%%
A second line of research originated with Stueckelberg and was developed
systematically by Horwitz, Piron and their collaborators.

The objective of this programme is not the construction of a stochastic
theory but rather the formulation of an unconstrained Hamiltonian
mechanics compatible with special relativity.
Its principal achievements include
\begin{itemize}
\item the introduction of an invariant evolution parameter,
\item manifestly covariant Hamiltonian dynamics,
\item off-shell particle dynamics,
\item a natural formulation of interacting many-particle systems.
\end{itemize}
The formalism provides an exceptionally attractive deterministic
framework.
However, stochastic dynamics does not form part of the theory in its
present form.
This programme is reviewed in detail in Chapter~2.
%%%%%%%%%%%%%%%%%%%%%%%%%%%%%%%%%%%%%%%%%%%%%%%%%%%%%%%%%%%%%%%%%%%%%%%%%%%%%%
\subsection{Relativistic Diffusion Processes}
%%%%%%%%%%%%%%%%%%%%%%%%%%%%%%%%%%%%%%%%%%%%%%%%%%%%%%%%%%%%%%%%%%%%%%%%%%%%%%
A third research programme attempts to construct stochastic processes
that are themselves compatible with Lorentz covariance.
Among the most important developments are relativistic Brownian motion,
diffusions on Lorentzian manifolds, stochastic processes on the mass
shell, and diffusion on the hyperboloid of four-velocities.
These investigations have established a mathematically rich theory of
relativistic stochastic processes.
Their relation to relativistic quantum mechanics, however, remains
incomplete.
These developments are reviewed in Chapter~4.
%%%%%%%%%%%%%%%%%%%%%%%%%%%%%%%%%%%%%%%%%%%%%%%%%%%%%%%%%%%%%%%%%%%%%%%%%%%%%%
\subsection{Relativistic Quantum Mechanics}
%%%%%%%%%%%%%%%%%%%%%%%%%%%%%%%%%%%%%%%%%%%%%%%%%%%%%%%%%%%%%%%%%%%%%%%%%%%%%%
The fourth programme is the conventional formulation of relativistic
quantum mechanics.
Beginning with the Klein--Gordon equation and culminating in the Dirac
equation and relativistic quantum field theory, this framework has
achieved extraordinary empirical success.
Its purpose differs fundamentally from that of stochastic mechanics.
Rather than deriving quantum theory from a deeper dynamical level, it
takes the wave equation as fundamental.
The relationship between these two viewpoints remains one of the
principal conceptual questions motivating the present work.
Chapter~5 reviews relativistic quantum mechanics from this perspective.
%%%%%%%%%%%%%%%%%%%%%%%%%%%%%%%%%%%%%%%%%%%%%%%%%%%%%%%%%%%%%%%%%%%%%%%%%%%%%%
\subsection{Current Status}
%%%%%%%%%%%%%%%%%%%%%%%%%%%%%%%%%%%%%%%%%%%%%%%%%%%%%%%%%%%%%%%%%%%%%%%%%%%%%%
Each of the four research programmes described above has solved
important parts of the overall problem.

Nelson's theory demonstrates that stochastic dynamics can reproduce
nonrelativistic quantum mechanics.

The SHP formalism provides a manifestly covariant Hamiltonian dynamics.

Relativistic diffusion theory establishes mathematically consistent
stochastic processes on Lorentzian manifolds.

Relativistic quantum mechanics provides the experimentally verified
description of elementary particles.

Nevertheless, none of these approaches simultaneously provides
\begin{enumerate}
\item manifest Lorentz covariance,
\item unconstrained Hamiltonian dynamics,
\item stochastic particle trajectories,
\item emergence of relativistic quantum mechanics.
\end{enumerate}
This observation defines the starting point of the present monograph.
%%%%%%%%%%%%%%%%%%%%%%%%%%%%%%%%%%%%%%%%%%%%%%%%%%%%%%%%%%%%%%%%%%%%%%%%%%%%%%
\subsection{Strategy of the Present Work}
%%%%%%%%%%%%%%%%%%%%%%%%%%%%%%%%%%%%%%%%%%%%%%%%%%%%%%%%%%%%%%%%%%%%%%%%%%%%%%
The approach adopted here is based upon the hypothesis that these
research programmes should not be regarded as competing alternatives.
Instead, they should be viewed as complementary developments, each
solving a different part of a larger problem.

The central objective of the present work is therefore to investigate
whether the deterministic Hamiltonian structure of the SHP formalism can
serve as the dynamical foundation upon which a relativistic stochastic
theory may be constructed.

Only after such a theory has been formulated will it become meaningful
to ask whether relativistic quantum mechanics emerges as its statistical
description.

